% Round 17. Root proof; independent review required.
% Reconstructs a historical R7 candidate through final KC17 and R16 inputs.
\section{Quadric forests and native ordinary nodes of arbitrary thickness}
\label{r17qf:section}

\paragraph{Provenance and scope.}
The node-chain geometry and a middle-degree energy calculation occur
in the historical Round 7 file
\path{proofs/good_ordinary_quadratic_higher_order.tex}, which remains
marked \texttt{NEEDS\_HUMAN\_REVIEW}.  Its nonmiddle proof invoked a
signed cohomological preimage interface which we do not import.
Here a numerical adjoint argument replaces that step.  The resulting
criterion applies to an arbitrary supplied forest with quadric double
strata, not just to chains with prescribed intermediate components.
The local geometric application is stated with exact pointed etale
equations, so no formal algebraization or characteristic-two normal
form is assumed.

All cycles and numerical groups have rational coefficients.  Let
$K/\mathbf Q_p$, $p\ne2$, have valuation ring $R$ and finite residue
field $k$.  Native degrees are always degrees over $k$; a quadric
defined over a finite field $k_e/k$ contributes the factor $[k_e:k]$.

\subsection{All-degree numerical correction on a quadric forest}

\begin{definition}[A native quadric forest model]
\label{r17qf:datum}
Let $\mathcal X/R$ be regular, integral, flat and projective,
strictly semistable of relative dimension $n\ge3$, with smooth
projective geometrically integral generic fiber.  Write its reduced
special fiber as $\bigcup_vD_v$.  Suppose its native incidence
multigraph, with a vertex for each native irreducible component and
an edge for each native connected double stratum, is a forest.
Suppose each double stratum $E_e$ is a smooth projective quadric of
dimension $n-1$ over its field of constants $k_e$, viewed as a
$k$-scheme.  The main components need not be quadrics or
geometrically integral over $k$.  Every closed stratum is smooth
over $k$.

A forest has no triple strata: three components through a point
would produce a triangle in the incidence multigraph.  Different
double strata between the same two components would produce parallel
edges and are excluded by the forest condition.
\end{definition}

Fix a relative ample line bundle $H$.  Write $H_v=H|_{D_v}$,
$h_e=H|_{E_e}$, and orient each edge.  In the notation of KC17 put
\[
 V^a=\bigoplus_vN^a(D_v)_{\mathbf Q},\qquad
 W^a=\bigoplus_eN^a(E_e)_{\mathbf Q},\qquad
 (\rho_a u)_e=u_{e^+}|_{E_e}-u_{e^-}|_{E_e}.
 \tag{QF17.1}\label{r17qf:rho}
\]
The correction map $A_a:V^a\to V^{a+1}$ is restriction of
$\sum_vi_{v*}u_v$.  For complementary edge degrees $a+b=n-1$,
the principal reduced fiber and projection formula give
\[
 \langle A_a u,w\rangle
       =-\sum_e\deg_k((\rho_a u)_e(\rho_b w)_e).
 \tag{QF17.2}\label{r17qf:adjoint}
\]
The component pairings here are unweighted native numerical
pairings.  The residue-field factors are already in $\deg_k$.

\begin{lemma}[Independent ample powers on the edges of a forest]
\label{r17qf:edge-powers}
For every $a\ge0$ and every edge $e$, the vector with value
$h_e^a$ at $e$ and zero at the other edges belongs to
$\operatorname{im}\rho_a$.
\end{lemma}

\begin{proof}
Remove $e$ from its tree, put $H_v^a$ on every vertex on one
side and zero on all other vertices.  Along each remaining edge
the two restrictions agree, because $H$ is one line bundle on
the entire model.  On $e$ their difference is $h_e^a$, up to
the orientation sign.  Changing the tuple's sign proves the
claim.  This includes $a=0$, with the actual component unit
classes.  No independent choices of polarizations on the
components are made.
\end{proof}

\begin{lemma}[A quadric middle form for an arbitrary ample class]
\label{r17qf:middle-form}
Let $E/k_0$ be a smooth projective quadric of dimension $2a$,
$a\ge1$, over a finite field of odd characteristic, and let
$h$ be any ample rational divisor.  On
\[
       P_h=\{u\in N^a(E)_{\mathbf Q}:\deg_{k_0}(u h^a)=0\},
\]
the form $(-1)^a\deg_{k_0}(u^2)$ is positive definite over
$\mathbf R$.
\end{lemma}

\begin{proof}
Pass to a finite splitting field only to compute the numerical
form.  Restriction is injective on native numerical groups,
and native degrees are preserved when each is measured over
its own base field.  If $a\ge2$, the geometric Picard group of
the quadric is generated by its natural hyperplane class $h_0$,
so $h$ is a positive rational multiple of $h_0$.  The two
middle ruling classes $l,l'$ have
$h_0^a=l+l'$ and
\[
       (l-l')^2=2(-1)^a
       \quad\hbox{under the degree map}.
\]
These are exactly the split intersection identities proved in
KC17, Lemma \texttt{r17kc:quadric}.  They prove the assertion
on the primitive line and hence on its native subspace.

If $a=1$, a split quadric is a quadric surface and
$h=\alpha l+\beta l'$ with $\alpha,\beta>0$.
Here $l^2=(l')^2=0$ and $\deg(ll')=1$.
The line orthogonal to $h$ is generated by
$\alpha l-\beta l'$ and has square $-2\alpha\beta<0$.
Again restrict to the native numerical subspace.  A nonsplit
quadric may have no nonzero native primitive vector; then
the assertion is vacuous.  No ordering of $\mathbf Q_\ell$
is used.
\end{proof}

\begin{theorem}[Exact numerical kernel equality on a quadric forest]
\label{r17qf:kernel}
For a model in Definition~\ref{r17qf:datum}, and every
$1\le r\le n$, put $a=n-r$.  Then
\[
                         \ker A_a=\ker\rho_a.
 \tag{QF17.3}\label{r17qf:kernel-equality}
\]
\end{theorem}

\begin{proof}
The inclusion $\ker\rho_a\subseteq\ker A_a$ follows at once
from \eqref{r17qf:adjoint} and perfect numerical duality on
the components.  For the other inclusion, take $u\in\ker A_a$,
put $\delta=\rho_a u$ and $b=r-1$, so $a+b=n-1$.

Suppose first that $a$ is not the middle degree of an
even-dimensional quadric.  Then neither is $b$.
For each $E_e$, both native numerical groups $N^a(E_e)$ and
$N^b(E_e)$ are one-dimensional ample lines.  This follows
from the split cellular groups, together with finite
restriction--corestriction and the nonzero native ample
class.  By Lemma~\ref{r17qf:edge-powers},
$\rho_b$ is onto the entire $W^b$.  Equation
\eqref{r17qf:adjoint} therefore makes $\delta$ orthogonal
to all of $W^b$.  The native numerical edge pairing is
perfect, so $\delta=0$.

The only other case is $n=2a+1$ and $r=a+1$, with $b=a$.
By \eqref{r17qf:adjoint} and
Lemma~\ref{r17qf:edge-powers}, each edge value satisfies
\[
                        \deg_k(\delta_e h_e^a)=0.
\]
Thus $\delta_e$ belongs to the primitive subspace of
Lemma~\ref{r17qf:middle-form}.  Pair the equation $A_a u=0$
with $u$ itself:
\[
                  0=-\sum_e\deg_k(\delta_e^2).
 \tag{QF17.4}\label{r17qf:energy}
\]
After multiplication by $(-1)^a$, the terms on the right
have one common definite sign.  More explicitly their
absolute sign-corrected values are
$[k_e:k](-1)^a\deg_{k_e}(\delta_e^2)\ge0$.
Every weight is positive, so all $\delta_e$ vanish.
There is no positivity assumption on any $D_v$.
\end{proof}

\begin{corollary}[All-degree native actual strong G]
\label{r17qf:actual}
On a quadric forest model every geometrically homologically
trivial rational class on the generic fiber has a numerically
admissible rational Chow extension in every positive codimension.
Each specified representative has a strong ordinary actual Green
form on its exact support complement, for every compact actual
test and the original valued-field normalization.
\end{corollary}

\begin{proof}
Numerically zero native cycles on each smooth quadric have
zero geometric cycle class in every degree, as proved in
KC17 by splitting and finite projection.  The exact equality
\eqref{r17qf:kernel-equality} therefore makes arbitrary actual
Chow representatives of the adjoint kernel cohomologically
compatible on every double stratum.  Apply the reduced case
of KC17's compatible-kernel criterion.  It produces an actual
rational vertical Chow correction.

Use the frozen pure-$K_0$ numerical witness, native all-degree
dominant detector, and the curved Green pair for the final
specified cycle.  These give obstruction zero and then the
claimed exact-support actual form.  The proof does not rely
on a signed cohomological boundary preimage, arbitrary proper
actual pushforward, or a finite-off-support assertion.
\end{proof}

\subsection{Exact native nodes of arbitrary positive order}

\begin{definition}[The supplied thick-node model]
\label{r17qf:thick-datum}
Let $\mathcal V/R$ be integral, flat and projective of relative
dimension $n\ge3$, with smooth geometrically integral generic
fiber $X$.  Suppose the special fiber is geometrically
integral and smooth except at finitely many closed points
$p_j$.  Let $k_j=\kappa(p_j)$ and let $R_j/R$ be the
corresponding finite unramified extension.  At each $p_j$
assume a residue-isomorphic pointed etale chart, with its
structural $R_j$-map supplied, to
\[
       \operatorname{Spec}
       R_j[x_0,\ldots,x_n]/(Q_j(x)-u_j\pi^{b_j}),
       \qquad b_j\ge1,\quad u_j\in R_j^\times,
 \tag{QF17.5}\label{r17qf:thick-equation}
\]
at the origin.  The quadratic form $Q_j$ is homogeneous and
has nondegenerate reduction.  The total model is not assumed
regular at these points.
\end{definition}

\begin{proposition}[Native semistable chains for the supplied equations]
\label{r17qf:thick-geometry}
For a model in Definition~\ref{r17qf:thick-datum}, make the
single extension $K'=K(\tau)$ with $\tau^2=\pi$.
At the node $p_j$, precisely $b_j$ successive blowups at the
unique remaining reduced point give a regular strictly
semistable model locally.  Globally the resulting model is
regular, projective and strictly semistable.  The fiber is
a tree consisting of the smooth strict transform of the
main component, and for each original point a chain with
$b_j-1$ intermediate components
\[
          \mathbf P_{E_j}
             (\mathcal O_{E_j}\oplus\mathcal O_{E_j}(-1)),
          \qquad
          E_j=\{\overline Q_j=0\}\subset\mathbf P^n_{k_j},
 \tag{QF17.6}\label{r17qf:intermediate}
\]
ending in the smooth quadric
\[
          T_j=\{\overline Q_j(X)=\overline u_j S^2\}
                           \subset\mathbf P^{n+1}_{k_j}.
 \tag{QF17.7}\label{r17qf:terminal}
\]
Every double stratum in that chain is a native copy of $E_j$,
every component has multiplicity one, and there are no triple
strata.  The two sections of each intermediate bundle are its
two disjoint neighboring double strata.
\end{proposition}

\begin{proof}
The extension is Eisenstein of full degree two and residue
degree one.  On the unramified chart base after this extension
the equation is
$Q(x)=u\tau^{2b}$.  Blow up the reduced maximal ideal
$(x_0,\ldots,x_n,\tau)$.  The saturated $x_i$-chart has
equations
\[
 x_i v=\tau,\qquad
 Q(y_0,\ldots,1,\ldots,y_n)
                          =u x_i^{2b-2}v^{2b};
 \tag{QF17.8}\label{r17qf:x-chart}
\]
the $\tau$-chart has equation
\[
                         Q(z)=u\tau^{2b-2}.
 \tag{QF17.9}\label{r17qf:tau-chart}
\]
These are actual Rees charts: removing the square of the
chart generator is saturation, and the resulting quotient
has no torsion by that generator.  For the $x_i$-chart this
can be checked by reducing the second equation modulo $x_i$:
its nonzero quadric polynomial is a nonzerodivisor in
$k_j[v,(y_l)]$.  The regular sequence with $x_i$ may be
permuted in the Cohen--Macaulay semistable ambient chart.
For the $\tau$-chart its defining polynomial is primitive
over the DVR, so multiplication by $\tau$ is injective.
Both charts thus embed in the corresponding localized Rees
chart and have exactly its saturation.

For $b>1$, at $x_i=0$ the reduced quadric equation in the
$y$ variables is smooth, as is its restriction at $v=0$.
Nondegeneracy and the identity $2Q=\sum x_i\partial_iQ$
show that one of the $y$ derivatives is nonzero on the
projective quadric chart.  Eliminating that variable etale
locally leaves $x_iv=\tau$.  Thus these charts are already
regular strictly semistable.
The $\tau$-chart has its only remaining nonsmooth special
point at $z=0$, with the order reduced from $2b$ to
$2b-2$.  The new exceptional component is the projective
cone over $E_j$; its vertex lies away from the preceding
double stratum.

The blowup of that cone at its vertex is the projective
bundle in \eqref{r17qf:intermediate}.  One explicit map sends
a line in $\mathcal O_{E_j}(-1)\oplus\mathcal O_{E_j}$ to
the corresponding point on the line from its base point
of $E_j$ to the cone vertex.  The vertex ideal pulls back
to the ideal of a section.  On each direction chart the
coordinates are $Y_l=Y_i y_l$, with
$\overline Q(y_0,\ldots,1,\ldots,y_n)=0$; these are precisely
the blowup charts.  The other distinguished section is the
unchanged hyperplane at infinity.  They are disjoint.

At $b=1$, the $\tau$-chart is $Q(z)=u$, smooth over the
DVR; the $x_i$-chart has $Q(1,y)=uv^2$.  Its exceptional
projective quadric is \eqref{r17qf:terminal}, smooth because
$2$ and $u$ are units.  The double locus $v=0$ is $E_j$.
This terminates the induction.  There are exactly $b$
blowups, $b-1$ intermediate bundles, and one terminal
quadric.  At each step the displayed structural equation
$x_iv=\tau$ proves multiplicity one.

The initial pointed etale maps identify the native completed
local rings and associated graded rings, so the exceptional
quadrics and cone bundles are defined over $k_j$ itself.
Blowup commutes with these etale maps.  Thus the regularity,
strict normal crossings, and native identifications descend.
The chains at different points are disjoint.  Away from the
points the original morphism is smooth by flatness and the
smoothness of its fiber; hence no other singularities or
components occur.  Blowups preserve projectivity, and
flatness and integrality follow from the displayed charts
and the integral generic fiber.

This is also the order-reduction calculation of
Jannsen--Saito, \emph{Bertini theorems and Lefschetz pencils
over discrete valuation rings}, Section 3, Theorem 5(i)
and its proof, printed pp.~12--13.  Their local normal-form
discussion is not used to enlarge the exact hypotheses
of Definition~\ref{r17qf:thick-datum}.
\end{proof}

\begin{theorem}[All-degree G for the supplied thick-node models]
\label{r17qf:thick-actual}
For the generic fiber in Definition~\ref{r17qf:thick-datum},
all geometrically homologically trivial rational cycles in
all positive codimensions satisfy strong ordinary actual G
over the original field $K$, on each specified exact support
complement and against every original compact actual test.
The same conclusion holds for a smooth projective target
having such a projective dominant witness after a finite
auxiliary local extension.
\end{theorem}

\begin{proof}
Proposition~\ref{r17qf:thick-geometry} gives a quadric forest
over $K'$.  Corollary~\ref{r17qf:actual} provides numerical
admissible lifts and actual obstruction zero there in every
degree.  The frozen native numerical-witness detector
descends obstruction zero to the original field, dividing
by the full generically finite and constant extension
degrees.  The auxiliary quadratic extension has degree
two, independently of the degrees of the node residue
fields.  Those fields enter only through the native
intersection degrees on the components.  Reconstruction
from the curved pair for the final prescribed cycle then
gives its exact-support Green form on the original space.
\end{proof}

\begin{corollary}[A complete height on a specified reduction class]
\label{r17qf:global}
Let $X/F$ be smooth projective geometrically integral over
a number field.  Assume that at every finite place, after
a finite auxiliary local extension if needed, $X$ has a
projective dominant witness with either good reduction or
a supplied thick-node model as in
Definition~\ref{r17qf:thick-datum}.  The thick-node alternative
is used only at places of odd residue characteristic.
Then the frozen PH16 height theorem supplies the complete
canonical real Beilinson--Bloch pairing on the geometric
homological rational Chow groups in every complementary
pair of degrees.
\end{corollary}

\begin{proof}
Good reduction gives the frozen numerical witness in every
degree.  The other local alternative is
Theorem~\ref{r17qf:thick-actual}.  Therefore all the local
strong-G hypotheses of PH16 hold in both complementary
degrees at every finite place.  PH16 gives rational
equivalence, symmetry, finite support of the sum, the
archimedean normalization, and full field-degree descent.
There is no additional unproved residue condition in this
application.  The reduction hypothesis is substantive and
has not been established for arbitrary number-field varieties.
\end{proof}

\paragraph{Remaining boundary.}
This argument does not cover an arbitrary graph with cycles:
the scalar ample edge vectors used in
Lemma~\ref{r17qf:edge-powers} need not span there.  Nor does
it prove characteristic-two ordinary normal forms, degenerate
quadrics, general singularities, or the existence of the
stated models for every variety.  Its historical components
are reactivated only through the exact proofs and final
actual interfaces given here.
